\subsection{Qué sí cubre la literatura revisada}

Esta subsección sintetiza, de forma transversal, lo que los catorce papers aportan al diseño del
modelo. No sustituye la lectura por ejes en §\ref{sec:estado-arte}.

\begin{enumerate}
  \item \textbf{OpenAPI} es el formato dominante para especificar APIs REST (E1). Casas et al.\
        (2021) mapean 47 artículos que lo confirman; Mainas et al.\ (2023) muestran enriquecimiento
        semántico viable sobre corpus reales.
  \item Existen marcos para \textbf{correspondencia entre esquemas} (E2) y \textbf{similitud
        semántica} (E3). Shvaiko y Euzenat (2005) y El Bayadh et al.\ (2015) ofrecen taxonomías
        reutilizables; Santos Sousa et al.\ (2026) documentan limitaciones prácticas de
        \textit{embeddings}.
  \item La \textbf{arquitectura orientada al dominio} separa modelos de negocio e interfaces
        técnicas (E4). Schmidt y Thiry (2020) relacionan capacidades de dominio con microservicios,
        en la misma línea que BIAN como referencia.
  \item Por el contrario, la literatura de \textbf{testing y seguridad} (E5) es abundante: Ehsan
        et al.\ (2022) y Golmohammadi et al.\ (2024) usan OpenAPI como entrada a pruebas
        automatizadas. Esa línea ayuda a \textbf{delimitar} el plan ---calidad operativa en
        producción no equivale a coherencia documental con BIAN.
  \item En selección de servicios (E6), Mahanra Rao et al.\ (2025) optimizan ranking por QoS con
        \textit{embeddings}; el objetivo es consumo, no verificar si un contrato bancario refleja un
        \ServiceDomain{} de referencia.
\end{enumerate}

Ninguna línea, sola o en conjunto, responde a la pregunta del plan: \textit{¿en qué medida un
contrato OpenAPI bancario refleja el modelo de un \ServiceDomain{} BIAN y cómo cuantificarlo de
forma reproducible?}

\subsection{Vacíos identificados}

\begin{table}[H]
  \centering
  \caption{Vacíos identificados en la literatura (V1--V5)}
  \label{tab:vacios}
  \small
  \begin{tabular}{@{}p{0.14\textwidth}p{0.38\textwidth}p{0.38\textwidth}@{}}
    \toprule
    \textbf{Vacío} & \textbf{Evidencia en el review} & \textbf{Implicación para la tesis} \\
    \midrule
    \textbf{V1 --- Intersección OpenAPI $\cap$ BIAN} &
    Ningún paper estudia ambos artefactos en el sector bancario &
    El objeto de estudio (contrato OpenAPI bancario vs.\ SD BIAN) no está cubierto \\
    \textbf{V2 --- Métrica integrada} &
    E2 y E3 aportan técnicas parciales; E1 no evalúa alineación &
    Justifica S, E, C y \AlignmentScore{} como producto original \\
    \textbf{V3 --- Cobertura de conceptos} &
    Matching ontológico discute cobertura; no sobre \BusinessObject{} BIAN en OpenAPI &
    Justifica \CoverageScore{} (C) \\
    \textbf{V4 --- Procedimiento reproducible} &
    Mainas propone \textit{pipeline} RDF; no scores comparables ni clasificación &
    Justifica diseño metodológico Cap.~4 (Guía~1) \\
    \textbf{V5 --- Gobernanza documental} &
    E5/E6 abordan \textit{runtime} o selección, no deriva OpenAPI--modelo de referencia &
    Motiva alineación previa a despliegue/certificación \\
    \bottomrule
  \end{tabular}
\end{table}


La \Cref{fig:brecha} sintetiza la transición desde los ejes del estado del arte hacia los vacíos
V1--V5 y la oportunidad de investigación.

\begin{figure}[H]
  \centering
  \includegraphics[width=0.95\textwidth]{Figura_06_Brecha}
  \caption{De la literatura revisada (E1--E6) a la brecha de investigación (V1--V5)}
  \label{fig:brecha}
\end{figure}

\subsection{Oportunidad de investigación}

La brecha abre la oportunidad de \textbf{diseñar un modelo de alineación} documentado en cinco
fases (Fase~1: diseño; Fases~2--5: transformaciones OE1--OE4). Dado un contrato OpenAPI bancario y
un \ServiceDomain{} BIAN emparejado, el modelo debe:

\begin{itemize}
  \item normalizar ambos artefactos a representaciones comparables (OE1);
  \item calcular \SemanticScore{}, \StructuralScore{} y \CoverageScore{} (OE2--OE3);
  \item integrar \AlignmentScore{} $= \alpha \cdot S + \beta \cdot E + \gamma \cdot C$ con
        clasificación interpretativa (Alta / Media / Baja / Nula).
\end{itemize}

Este plan documenta el \textbf{diseño} del procedimiento. La construcción del artefacto y la
experimentación empírica quedan fuera de alcance (\Cref{sec:alcances}).

La investigación contribuye a los vacíos V1--V4 con un enfoque reproducible sobre artefactos
documentales. No sustituye testing operativo (E5) ni certificación regulatoria; tampoco confunde
calidad \textit{runtime} (E1/E6) con coherencia documental (V5).
